\subsection*{NT-A06--A07. Gaussian prime angles at almost every fixed center}

Write \(\mathbb T_G=\mathbf R/(\pi/2)\mathbf Z\), with probability
measure \(d\mu_G(\theta)=2\,d\theta/\pi\). For a nonzero ideal
\(\mathfrak a=(\alpha)\) of \(\mathbf Z[i]\), its angle is
\(\theta_{\mathfrak a}=\arg\alpha\pmod{\pi/2}\). Put
\[
 \Xi_m((\alpha))=(\alpha/\overline\alpha)^{2m},
 \qquad m\in\mathbf Z .
\]
These are well-defined conductor-one unitary Hecke characters.
For \(m\ne0\) their finite-part functions are entire, with completion
\[
 \pi^{-(s+2|m|)}\Gamma(s+2|m|)L(s,\Xi_m)
 =
 \pi^{-(1-s+2|m|)}\Gamma(1-s+2|m|)L(1-s,\Xi_m).
\]
The precise analytic assumption in this subsection is RH for
\(L(s,\Xi_m)\) for every positive integer \(m\); negative \(m\)
give the same functions. Conditional input \(M\) supplies this family.
The mean prime count below uses only the unconditional prime ideal
theorem for \(\mathbf Q(i)\). Input \(A\) alone does not supply the
nonzero angular infinity types.

\paragraph{Classical variance input.}
Rudnick--Waxman, \emph{Angles of Gaussian primes},
\texttt{arXiv:1705.07498v2}, Theorem~3.2, author p.~12,
proved in \S4.4, p.~20, gives the following. Fix nonzero nonnegative
functions \(f\in C_c^\infty((-1/2,1/2))\), even, and
\(\Phi\in C_c^\infty((1/2,1))\). Define
\[
 F_K(\theta)=\sum_{j\in\mathbf Z}
 f\left(\frac{K}{\pi/2}(\theta-j\pi/2)\right),\qquad
 P_{K,X}(\theta)=
 \sum_{\mathfrak p}\Phi(N\mathfrak p/X)\log N\mathfrak p\,
                         F_K(\theta_{\mathfrak p}-\theta).
\]
In the range used here, \(3\le K\le X\le K^2\), that theorem gives
\[
 \operatorname{Var}_{\mu_G}(P_{K,X})
       \ll_{f,\Phi}\frac{X}{K}(\log K)^2.             \tag{A06.1}
\]
Lemma~3.1, author pp.~11--12, gives
\[
 \mathbb E_{\mu_G}P_{K,X}
       \sim c_{f,\Phi}\frac XK,\qquad
 c_{f,\Phi}=\left(\int_{\mathbf R}f\right)
            \left(\int_0^\infty\Phi\right)>0.        \tag{A06.2}
\]
The asymptotic is uniform in \(K\): its entire \(K\)-dependence is
the exact factor \(1/K\) in the integral of \(F_K\).
The nonzero Fourier modes in \S4 use precisely the functions
\(L(s,\Xi_m)\) above; no zero-correlation or simplicity hypothesis
is part of the variance upper bound. The full asymptotic variance
conjecture of that paper is not used.

\paragraph{NT-A06: an explicit exceptional-measure formulation.}
Let \(J_K(\theta)\) be the open centered arc of length \(\pi/(2K)\).
For \(X\) sufficiently large and \(3\le K\le X\le K^2\),
let \(E_{K,X}\) be the set of centers whose arc \(J_K(\theta)\)
contains no prime ideal angle of norm in \((X/2,X)\), and let
\(E^{(1)}_{K,X}\) impose the same absence for absolute-degree-one
prime ideals. Then
\[
 \mu_G(E_{K,X})\ll\frac{K(\log K)^2}{X},\qquad
 \mu_G(E^{(1)}_{K,X})\ll\frac{K(\log K)^2}{X}+\frac1K .
                                                               \tag{A06.3}
\]
The implied constants are fixed once the two auxiliary functions
are fixed; no numerical constant is claimed.

\begin{proof}
For all sufficiently large \(X\), (A06.2) gives
\(\mathbb E P_{K,X}\ge c_{f,\Phi}X/(2K)>0\), uniformly in \(K\).
If \(E_{K,X}\) occurs, every summand in \(P_{K,X}\)
vanishes by its support. Chebyshev's inequality therefore gives
\[
 \mu_G\{P_{K,X}=0\}
 \le \frac{\operatorname{Var}(P_{K,X})}
           {(\mathbb E P_{K,X})^2}
 \ll \frac{K(\log K)^2}{X}.
 \]
The only prime ideals of absolute degree greater than one in
\(\mathbf Q(i)\) are those above inert rational primes, all with
angle zero. The set of centers whose arc \(J_K(\theta)\) contains
zero has measure \(1/K\). Outside that set any prime ideal found
by the first assertion has degree one. This proves the second
assertion, including the endpoint conventions (endpoints form a
measure-zero set).
\end{proof}

\paragraph{NT-A07: almost every center, at every sufficiently fine scale.}
Fix \(\epsilon>0\). Under the same angular RH family, for
\(\mu_G\)-almost every fixed \(\theta\), there is
\(K_0(\theta,\epsilon)\) such that for every real \(K\ge K_0\),
\(J_K(\theta)\) contains the angle of a degree-one Gaussian prime
ideal satisfying
\[
 N\mathfrak p\le C_\epsilon
              K(\log K)^3(\log\log K)^{1+\epsilon}.       \tag{A07.1}
\]
The constant \(C_\epsilon\) is independent of the center.
For almost every center this produces infinitely many such prime
ideals. The threshold may depend on the center. This is a
complete deduction from the cited variance bound, not a claim
of priority over other almost-sure formulations.

\begin{proof}
For integers \(j\) sufficiently large put
\[
 K_j=2^j,\qquad
 X_j=K_j(\log K_j)^3(\log\log K_j)^{1+\epsilon}.
\]
Then \(K_j\le X_j\le K_j^2\). Let \(E_j\) be the event that
no degree-one prime ideal of norm in \((X_j/2,X_j)\) has angle
in \(J_{K_j}(\theta)\). Formula (A06.3) gives
\[
 \mu_G(E_j)\ll_\epsilon
      \frac1{j(\log j)^{1+\epsilon}}+2^{-j}.
\]
This series converges. The first Borel--Cantelli lemma, which
requires no independence, says that almost every center belongs
to only finitely many \(E_j\). Consequently every sufficiently
large dyadic scale has such a prime.

For arbitrary sufficiently large real \(K\), choose \(j\) with
\(K\le 2^j<2K\). The centered arc \(J_{K_j}(\theta)\) is contained
in \(J_K(\theta)\). For all sufficiently large \(K\),
\(\log K_j\le2\log K\) and
\(\log\log K_j\le2\log\log K\); hence
\[
 X_j\le 2^{5+\epsilon}
           K(\log K)^3(\log\log K)^{1+\epsilon}.
\]
This proves (A07.1), for example with \(C_\epsilon=2^{5+\epsilon}\)
after increasing the threshold. This constant comes only from
the dyadic comparison, not from an explicit numerical variance
estimate. The selected dyadic primes have norm at least \(X_j/2\),
which tends to infinity, so infinitely many distinct prime ideals
occur.
\end{proof}

\paragraph{Quantifier boundary and geometric placement.}
Almost every fixed center cannot be replaced by a bound uniform
over all centers at this norm scale. If a nonreal Gaussian
integer \(a+ib\) has \(a,b>0\) and norm at most \(X\), then
\(\arg(a+ib)\ge\arctan(1/\sqrt X)\). Hence the open arc
\((0,\arctan(1/\sqrt X))\) contains no degree-one prime angle
of norm at most \(X\). A Gaussian prime may be represented in
the closed first quadrant after multiplication by a unit;
the degree-one primes on the axes do not give angles in that
open arc. This elementary gap explains why the exceptional set
and its order of quantifiers matter.

In the Grothendieck/Frobenius dictionary the prime ideals are
closed points of \(\operatorname{Spec}\mathbf Z[i]\), and the
\(\Xi_m\) are the integer weights of its compact CM torus.
The variance averages over the center of a torus neighborhood;
the first Borel--Cantelli step turns that averaged bound into
a statement about almost every fixed neighborhood system.
No independent random primes or independent zeros are assumed.
Run location: \texttt{GRH\_NIGHT/number\_theory};
the exact source is \texttt{S-NT-RW}, and the full new deductions
are (A06.3) and NT-A07 above.
